\documentclass[12pt]{article}
\usepackage[utf8]{inputenc}
\usepackage[T1]{fontenc}
\usepackage[margin=1in]{geometry}
\usepackage{setspace}
\usepackage{xcolor}
\usepackage[colorlinks=true, urlcolor=blue!60!black, linkcolor=blue!60!black]{hyperref}
\setlength{\parindent}{2em}
\setlength{\parskip}{1.25em}
\renewcommand{\baselinestretch}{1.0}
\begin{document}
\begin{center}
{\LARGE \textit{U.S. Term Limits, Inc. v. Thornton}, 514 U.S. 779 (1995)\par}
\vspace{0.8em}
{\large Case Brief\par}
\end{center}
\vspace{1em}

\section{Memory Jogger}
Arkansas barred three-term Representatives and two-term Senators from its general-election ballot; a 5--4 Court held that the Constitution's qualifications for Congress are fixed and exclusive, so neither Congress nor a State may add term limits, whether directly or disguised as ballot-access rules---only an Article V amendment can (\href{https://supreme.justia.com/cases/federal/us/514/779/}{Justia}).

\section{Detailed Case Facts}
On November 3, 1992, Arkansas voters approved Amendment 73 to the Arkansas Constitution, a sweeping term-limits initiative whose preamble declared that ``[e]ntrenched incumbency'' had made elections ``less free, less competitive, and less representative.'' Sections 1 and 2 imposed direct term limits on Arkansas's state executive and legislative officials. Section 3 addressed the Arkansas congressional delegation: rather than formally barring long-serving incumbents from holding federal office, it prohibited any person who had been elected to three or more terms in the U.S. House of Representatives or two or more terms in the U.S. Senate from being certified as a candidate or having his or her name printed on the general-election ballot, leaving only the option of a write-in campaign.

\section{Procedural History}
Ten days later, on November 13, 1992, Arkansas voter Bobbie Hill---suing on behalf of herself, similarly situated Arkansas voters, and the League of Women Voters of Arkansas---filed a declaratory-judgment action in Pulaski County Circuit Court challenging \S~3 under the Federal Constitution. Arkansas Attorney General Winston Bryant and proponents including U.S. Term Limits, Inc. intervened to defend the amendment, while U.S. Representative Ray Thornton joined on the challengers' side. The Circuit Court held \S~3 unconstitutional under Article I. On appeal, a 5--2 Arkansas Supreme Court affirmed, with a plurality concluding that States lack authority to alter the Qualifications Clauses and rejecting Arkansas's characterization of \S~3 as a mere ``ballot access'' rule. By the time the U.S. Supreme Court granted certiorari, 22 States had enacted congressional term-limit or ballot-access measures.

\section{Judicial Votes}
The Supreme Court \textbf{affirmed} the Supreme Court of Arkansas by a \textbf{5--4} vote. The case was argued on November 29, 1994, and decided on May 22, 1995 (\href{https://supreme.justia.com/cases/federal/us/514/779/}{Justia}; \href{https://www.law.cornell.edu/supremecourt/text/514/779}{Cornell LII}):

\begin{itemize}
    \item \textbf{Majority Opinion:} Justice Stevens delivered the opinion of the Court, joined by Justices Kennedy, Souter, Ginsburg, and Breyer.
    \item \textbf{Concurring Opinion:} Justice Kennedy filed a concurring opinion.
    \item \textbf{Dissenting Opinion:} Justice Thomas filed a dissenting opinion, joined by Chief Justice Rehnquist and Justices O'Connor and Scalia.
\end{itemize}

\section{Holding}
\textbf{Questions Presented:}

\begin{enumerate}
    \item Do the Qualifications Clauses of Article I (\S~2, cl. 2 and \S~3, cl. 3) establish fixed and exclusive qualifications for service in Congress, or do States retain reserved power under the Tenth Amendment to add congressional qualifications such as term limits?
    \item Does the Elections Clause (U.S. Const. art. I, \S~4, cl. 1) empower a State to bar multi-term congressional incumbents from the general-election ballot as a ``Manner'' regulation while leaving them theoretically eligible to win via write-in votes?
\end{enumerate}

\textbf{Holding:} \textbf{Affirmed} (5--4); \S~3 of Arkansas Amendment 73 is \textbf{unconstitutional}. The Qualifications Clauses of Article I are fixed and exclusive: neither Congress nor the individual States may add to, alter, or diminish the age, citizenship, and inhabitancy requirements enumerated in the Constitution. States possess no reserved Tenth Amendment authority over federal congressional offices, nor may they use the Elections Clause to impose indirect term limits through ballot-access restrictions designed to handicap a class of candidates.

\section{Analysis of Opinions}
\subsection{Majority Opinion (Justice Stevens)}

Writing for a five-Justice majority, Justice Stevens framed the case around the nature of national sovereignty and the democratic principles recognized in \href{https://supreme.justia.com/cases/federal/us/395/486/}{\textit{Powell v. McCormack}, 395 U.S. 486 (1969)}.

\subsubsection{Reaffirming \textit{Powell} and Rejecting Reserved State Power}

First, Stevens reaffirmed \textit{Powell}'s holding that Congress itself cannot add qualifications beyond Article I's text. Turning to state power, Stevens rejected petitioners' Tenth Amendment defense on two independent grounds. First, adopting Justice Joseph Story's \textit{Commentaries} and Chief Justice Marshall's reasoning in \textit{McCulloch v. Maryland}, 17 U.S. (4 Wheat.) 316 (1819), Stevens held that the Tenth Amendment can ``reserve'' only pre-existing sovereign powers. Because the federal Congress did not exist before ratification of the 1787 Constitution, electing federal representatives was a new right arising from the Constitution itself; States cannot reserve a power they never possessed. Second, even assuming States held some original power, the constitutional text, structure, and history demonstrate that the Framers divested States of it. Members of Congress are paid from the national treasury (art. I, \S~6), their qualifications are judged by a national body (art. I, \S~5, cl. 1), and the Framers feared state interference with the National Legislature. Moreover, during the heated ratification debates over ``rotation in office,'' no Federalist ever suggested States could unilaterally enact congressional term limits, and immediately after ratification---while many States imposed property and term-limit qualifications on state officials---virtually none did so for Congress.

\subsubsection{Ballot-Access Disguise and the Elections Clause}

Second, Stevens rejected Arkansas's argument that \S~3 merely regulated ballot access under the Elections Clause (art. I, \S~4, cl. 1) because barred incumbents could still run as write-in candidates. A State cannot accomplish indirectly what the Constitution forbids directly. Because \S~3 had the ``likely effect of handicapping a class of candidates'' and the ``sole purpose of creating additional qualifications indirectly,'' it violated the Qualifications Clauses regardless of whether a write-in victory was theoretically possible. Reading the Elections Clause's ``Times, Places and Manner'' language to authorize substantive exclusions would also produce an absurd consequence under \textit{Powell}: Congress's power to ``make or alter'' state manner regulations would empower Congress itself to impose qualifications. Rather, the Elections Clause authorizes only procedural regulations ensuring fair, orderly elections (such as the anti-raiding and petition rules upheld in \textit{Storer v. Brown}, 415 U.S. 724 (1974)), not measures designed to dictate electoral outcomes. Any fundamental change such as congressional term limits must come through an Article V constitutional amendment.

\subsection{Concurring Opinion (Justice Kennedy)}

Providing the decisive fifth vote, Justice Kennedy wrote separately to articulate a foundational theory of dual sovereignty. He rejected the dissent's premise that the Constitution rests solely on the consent of the people of each separate State, emphasizing that the Framers ``split the atom of sovereignty'' so that every American holds two distinct political capacities---one state and one federal---each protected from intrusion by the other. Just as the Federal Government may not invade state sovereignty (citing \textit{United States v. Lopez}, 514 U.S. 549 (1995), decided one month earlier), a State may not interpose itself between national citizens and their National Government by burdening the federal right to vote for congressional representatives.

\subsection{Dissenting Opinion (Justice Thomas)}

In a sweeping 80-page dissent joined by Chief Justice Rehnquist and Justices O'Connor and Scalia, Justice Thomas advanced a fundamentally different vision of the Founding. Because the Constitution was ratified state by state under Article VII, Thomas argued that its ultimate authority derives from ``the people of each individual State,'' not an undifferentiated national people. Under the Tenth Amendment's default rule, where the Constitution is silent, the Federal Government lacks power and the States (or their people) retain it. Story's rule that States cannot reserve powers over newly created federal offices contradicts \textit{McCulloch} and cannot explain why States undisputedly set qualifications for presidential electors under Article II.

On the text, Thomas argued that the Qualifications Clauses establish only minimum \textit{floors} (preventing States from sending immature or non-citizen legislators), not exclusive \textit{ceilings}; the express prohibition on state religious tests in Article VI, cl. 3 would be superfluous if States were already barred from adding qualifications. Thomas further pointed to five early States (Virginia, Maryland, Massachusetts, Georgia, and North Carolina) that imposed district-residency or freehold requirements in the first federal elections, and he read \textit{Powell} narrowly as holding only that Article I, \S~5 did not grant \textit{Congress} power to exclude members. Finally, Thomas argued that \S~3 did not bar incumbents from serving---citing successful write-in campaigns---and could be viewed as a legitimate effort to offset federal incumbency advantages (such as franking privileges and seniority rules).

\subsection{Subsequent Developments}

\begin{itemize}
    \item \textbf{Anti-Evasion and Ballot Labels (\textit{Cook v. Gralike}):} In \href{https://supreme.justia.com/cases/federal/us/531/510/}{\textit{Cook v. Gralike}, 531 U.S. 510 (2001)}, the Court struck down a Missouri constitutional amendment instructing congressional candidates to support a federal term-limits amendment and printing ``scarlet letter'' ballot labels (``DISREGARDED VOTERS' INSTRUCTION ON TERM LIMITS'') next to noncompliant candidates. Relying directly on \textit{Thornton}, the Court held that States possess no reserved power over congressional elections and cannot misuse the Elections Clause to ``dictate electoral outcomes.''
    \item \textbf{Contrast with Presidential Electors (\textit{Chiafalo v. Washington}):} In \href{https://supreme.justia.com/cases/federal/us/591/19-465/}{\textit{Chiafalo v. Washington}, 591 U.S. 578 (2020)}, a unanimous Court upheld Washington's authority to fine ``faithless'' presidential electors. Crucially, the Court harmonized \textit{Chiafalo} with \textit{Thornton} by emphasizing that Article II, \S~1, cl. 2 expressly delegates broad power to each State to appoint electors ``in such Manner as the Legislature thereof may direct,'' whereas Article I delegates no comparable substantive authority over congressional candidates. (Justice Thomas concurred only in the judgment, reiterating his \textit{Thornton} dissent's Tenth Amendment theory.)
    \item \textbf{Federal Ballot-Access Disqualifications (\textit{Trump v. Anderson}):} In \href{https://supreme.justia.com/cases/federal/us/601/23-719/}{\textit{Trump v. Anderson}, 601 U.S. 100 (2024)} (per curiam), a unanimous Court reversed the Colorado Supreme Court's order barring former President Donald Trump from the 2024 presidential primary ballot under Section 3 of the Fourteenth Amendment. Citing \textit{Thornton} repeatedly, the Court held that States lack authority under the Elections or Electors Clauses to enforce Section 3 against federal candidates because state-by-state enforcement would create a chaotic ``patchwork'' that could ``undermine the National Government.''
\end{itemize}

\section{Examples: Future Applications}
\begin{enumerate}

\item \textbf{Same-Side (Unconstitutional): District Residency and a Tax-Return Ballot Label.} Seeking to combat corruption and promote local accountability, State X amends its state constitution to (1) require every U.S. House candidate to reside inside the specific congressional district for at least two years before the election, and (2) print ``REFUSED TAX-RETURN DISCLOSURE'' beside the name of any congressional candidate who declines to release five years of federal tax returns. \textit{Analysis:} Both provisions are \textbf{unconstitutional} under \textit{Thornton} and \href{https://supreme.justia.com/cases/federal/us/531/510/}{\textit{Cook v. Gralike}, 531 U.S. 510 (2001)}. Article I, \S~2, cl. 2 requires only that a Representative be an inhabitant of the \textit{State} ``when elected''; State X cannot add a durational district-residency qualification. Likewise, printing a pejorative ballot label to coerce tax disclosure is not a neutral ``Manner'' regulation under Article I, \S~4, cl. 1, but an unconstitutional effort to handicap a class of federal candidates and dictate electoral outcomes.

\item \textbf{Same-Side (Unconstitutional): Congress Limits Its Own Members' Terms by Statute.} Responding to public pressure, Congress passes a statute providing that no person who has served six terms in the House or two terms in the Senate is eligible for election to that chamber again. A long-serving Representative sues to stay on the ballot. \textit{Thornton} reaffirmed \textit{Powell v.\ McCormack}'s holding that the Qualifications Clauses are exclusive and that Congress itself cannot add to them, and Stevens's opinion explained that any change to those qualifications must come through Article V. A statutory term limit is an added qualification for congressional service, so it is unconstitutional no matter which body enacts it (\href{https://supreme.justia.com/cases/federal/us/514/779/}{Justia}).

\item \textbf{Opposite-Side (Constitutional): State Term Limits, Resign-to-Run, and Neutral Ballot Rules.} Now suppose State X instead (1) imposes a two-term limit on \textit{State X's own Governor and state legislators}, (2) enforces a ``resign-to-run'' rule requiring sitting state judges to resign their state office upon declaring candidacy for Congress, and (3) enacts a uniform ballot-access rule requiring any independent congressional candidate to submit signatures from 2 percent of registered voters in the district by June 1 and prohibiting sore-loser candidates who lost a party primary from appearing on the general-election ballot. The first two measures do not add qualifications for Congress. The third does not do so either, but its constitutionality also depends on the burden it imposes on ballot access. Term limits on \textit{state} officials rest squarely within State X's reserved Tenth Amendment sovereignty; ``resign-to-run'' laws permissibly regulate \textit{state officeholders} rather than imposing a qualification on federal office (which Stevens expressly distinguished in \textit{Thornton}, 514 U.S. at 835 n. 48, citing \textit{Clements v. Fashing}, 457 U.S. 957 (1982)); and evenhanded signature thresholds and sore-loser rules can be valid Elections Clause regulations. A court would assess the combined effect of the June 1 deadline, the signature requirement, and the State's other ballot-access rules against the State's election-administration interests.

\item \textbf{Opposite-Side (Constitutional): A State Applies for an Article V Convention.} Frustrated by \textit{Thornton}, a state legislature adopts a resolution applying to Congress, under Article V, to call a convention for proposing a constitutional amendment imposing congressional term limits, and urges its congressional delegation to support such an amendment. Opponents argue that the State is again trying to impose term limits it may not adopt. Unlike Arkansas's Amendment 73, the resolution adds no qualification, keeps no candidate off the ballot, and does not handicap any class of candidates. It pursues precisely the route \textit{Thornton} identified as the only legitimate path to congressional term limits: amendment of the Constitution under Article V, in which the Constitution gives state legislatures an express role.

\item \textbf{Fence-Sitter (Neutral Information or a Handicap?): A Ballot Notation of Terms Served.} A State requires its ballots to list, beside each congressional incumbent's name, the number of terms the incumbent has served (for example, noting that a Representative is completing a sixth term). The notation is accurate and applies to every incumbent, and the State says it simply informs voters. But the measure was adopted by an initiative whose sponsors openly sought to promote congressional term limits, and the label may be designed to disadvantage long-serving incumbents. Under \textit{Thornton} and \textit{Cook v.\ Gralike}, a ballot rule is invalid if it is meant to handicap a class of candidates and dictate electoral outcomes, but neutral, factual ballot information is a familiar Elections Clause regulation. Whether this notation is the former or the latter is genuinely uncertain.

\end{enumerate}

\section{Critique}
\textbf{Competing Founding Theories:} \textit{Thornton} exposes the deepest unresolved philosophical divide in constitutional federalism---whether the Constitution was created by one national people (Stevens and Kennedy) or by the separate peoples of sovereign States (Thomas). Neither the Preamble nor Article VII conclusively resolves that tension.

\textbf{Asymmetry of Incumbency:} As Thomas emphasized, \textit{Thornton} bars States from offsetting structural advantages that federal incumbents grant themselves (taxpayer-funded staffs, franking, seniority), while leaving Congress with little political incentive to propose an Article V term-limits amendment against its own members.

\textbf{Line-Drawing Under the Elections Clause:} Distinguishing permissible ``procedural'' ballot-access rules (such as \textit{Storer}'s disaffiliation requirement) from impermissible ``substantive'' qualifications turns heavily on judicial assessments of legislative purpose and whether a rule ``handicaps a class of candidates.''

\section{Key Quotations}
All quotations are verified verbatim from Cornell LII (\href{https://www.law.cornell.edu/supct/html/93-1456.ZO.html}{Majority}, \href{https://www.law.cornell.edu/supct/html/93-1456.ZC.html}{Concurrence}, \href{https://www.law.cornell.edu/supct/html/93-1456.ZD.html}{Dissent}) and \href{https://supreme.justia.com/cases/federal/us/514/779/}{Justia}:

\begin{quote}
``Contrary to petitioners' assertions, the power to add qualifications is not part of the original powers of sovereignty that the Tenth Amendment reserved to the States. Petitioners' Tenth Amendment argument misconceives the nature of the right at issue because that Amendment could only `reserve' that which existed before. As Justice Story recognized, `the states can exercise no powers whatsoever, which exclusively spring out of the existence of the national government, which the constitution does not delegate to them. \ldots{} No state can say, that it has reserved, what it never possessed.' ''
\end{quote}
\noindent --- Justice Stevens, 514 U.S. at 802 (\href{https://www.law.cornell.edu/supct/html/93-1456.ZO.html}{Cornell LII}).

\begin{quote}
``We are, however, firmly convinced that allowing the several States to adopt term limits for congressional service would effect a fundamental change in the constitutional framework. Any such change must come not by legislation adopted either by Congress or by an individual State, but rather---as have other important changes in the electoral process---through the Amendment procedures set forth in Article V.''
\end{quote}
\noindent --- Justice Stevens, 514 U.S. at 837 (\href{https://www.law.cornell.edu/supct/html/93-1456.ZO.html}{Cornell LII}).

\begin{quote}
``Federalism was our Nation's own discovery. The Framers split the atom of sovereignty. It was the genius of their idea that our citizens would have two political capacities, one state and one federal, each protected from incursion by the other.''
\end{quote}
\noindent --- Justice Kennedy, concurring, 514 U.S. at 838 (\href{https://www.law.cornell.edu/supct/html/93-1456.ZC.html}{Cornell LII}).

\begin{quote}
``Nothing in the Constitution deprives the people of each State of the power to prescribe eligibility requirements for the candidates who seek to represent them in Congress. The Constitution is simply silent on this question. And where the Constitution is silent, it raises no bar to action by the States or the people.''
\end{quote}
\noindent --- Justice Thomas, dissenting, 514 U.S. at 845 (\href{https://www.law.cornell.edu/supct/html/93-1456.ZD.html}{Cornell LII}).


\end{document}
